\begin{abstract}
Simultaneous recognition of facial demographic attributes such as race and gender faces challenges from expression variations, biological aging, phenotypic overlap, and the limitations of single-domain representations. This study proposes a multi-domain latent feature fusion framework that integrates visual representations from three task-specific pre-trained Vision Transformer (ViT) models for facial biometrics, affective expression, and age estimation, coupled with classical machine learning pipeline optimization for classifying six intersectional demographic subgroups on the DemogPairs dataset. Latent representations are extracted offline from each model to produce a unified visual representation that combines complementary cross-domain information. Seven feature configurations are evaluated through 5-Fold Stratified Cross-Validation across four classifiers: Random Forest (RF), Gaussian Naive Bayes (GNB), Logistic Regression (LR), and Support Vector Machine (SVM) with hyperparameter tuning via Grid Search Cross-Validation (GridSearchCV). Experimental results demonstrate that tri-domain fusion achieves the best performance across three of the four evaluated classifiers, with the SVM model yielding the highest performance, achieving an Accuracy of 93.70\%, Precision of 93.72\%, Recall of 93.70\%, and F1-Score of 93.69\% on independent test data. In-depth evaluation across the six subgroups records an F1-Score range between 91.74\% and 96.14\%, proving the effectiveness of multi-domain latent representation fusion in recognizing intersectional demographic attributes.
\end{abstract}

\begin{keywords}
Race and gender classification, intersectional demographic recognition, Vision Transformer, multi-domain feature fusion, algorithmic fairness.
\end{keywords}

\titlepgskip=-21pt
